\section*{Bab \ref{ch:b3:chromatin-epigenetics} --- Kromatin dan Epigenetika}

\begin{solution}{exo:b3:chromatin-epigenetics:1}
Masing-masing dua salinan H2A, H2B, H3, dan H4 (sebuah oktamer) dengan
\qty{147}{bp} DNA dalam $1.65$ putaran; penghubung sepanjang
\qtyrange{20}{60}{bp}, sehingga ulangannya mendekati \qty{200}{bp}.
Nukleasenya hanya memotong DNA penghubung yang terpapar, tak pernah di dalam
sebuah inti, sehingga setiap fragmen merupakan kelipatan bulat ulangan itu:
sebuah tangga \num{200}, \num{400}, \qty{600}{bp}, bukan lapisan merata.
\end{solution}

\begin{solution}{exo:b3:chromatin-epigenetics:2}
Diploid: $6.0\times 10^{9}$ bp; $6.0\times 10^{9}/190 \approx
3.2\times 10^{7}$ \omterm{def:b3:chromatin-epigenetics:nucleosome}{nukleosom}; panjangnya $6.0\times 10^{9}\times
\qty{0.34}{nm} = \qty{2.0}{m}$.
\end{solution}

\begin{solution}{exo:b3:chromatin-epigenetics:3}
H3K27me3: \omterm{def:b3:chromatin-epigenetics:nucleosome}{histon} H3, lisina 27, termetilasi tiga kali --- \omterm{def:b3:chromatin-epigenetics:nucleosome}{kromatin} bisu gen
yang ditekan selama perkembangan, dituliskan PRC2 dan dibaca protein
berkromodomain PRC1. H4K16ac: H4, lisina 16, terasetilasi --- \omterm{def:b3:chromatin-epigenetics:nucleosome}{kromatin} aktif
dan terbuka (tanda itu memutus satu sentuhan antarnukleosom). H3S10ph: H3,
serina 10, terfosforilasi --- tanda mitosis yang berkaitan dengan pemadatan
kromosom (dan, sesaat, dengan pengaktifan gen oleh sebagian isyarat); bukan
tanda pembungkaman maupun pengaktifan yang mantap.
\end{solution}

\begin{solution}{exo:b3:chromatin-epigenetics:4}
Gen jingga terpaut-X dan memiliki alel jingga serta alel hitam; bulu
berbercak menuntut satu alel dari masing-masing, pada dua kromosom X yang
berbeda, dan salah satunya dibungkam di tiap sel. Jantan memiliki satu X dan
seluruhnya jingga atau seluruhnya hitam. Jantan belang tiga membawa dua
kromosom X dan sebuah Y (XXY), dan mandul.
\end{solution}

\begin{solution}{exo:b3:chromatin-epigenetics:5}
$\qty{200}{bp}\times \qty{0.34}{nm} = \qty{68}{nm}$ tiap manik setebal
\qty{11}{nm}: rasionya $6.2$. Yang dibutuhkan: $\qty{4}{cm}/\qty{4}{\micro m} =
10^{4}$. Pelipatan lanjutannya sebesar $10^{4}/6.2 \approx 1600$.
\end{solution}

\begin{solution}{exo:b3:chromatin-epigenetics:6}
$\mu = 0.98$, $\delta = 0.01$: $p^{*} = 0.01/0.03 = 0.33$; rasionya $\mu -
\delta = 0.97$; waktu paruhnya $\ln 2/(-\ln 0.97) = 22.8 \approx 23$
pembelahan. $\mu = 0.999$, $\delta = 0.001$: $p^{*} = 0.5$; rasionya
$0.998$; waktu paruhnya $\ln 2 / 0.002 \approx 350$ pembelahan.
\end{solution}

\begin{solution}{exo:b3:chromatin-epigenetics:7}
$A$: aktif (promotor H3K4me3 dengan \omterm{def:b3:chromatin-epigenetics:histone-code}{asetilasi}). $B$: bisu, ditekan
Polycomb --- \omterm{def:b3:chromatin-epigenetics:nucleosome}{heterokromatin} fakultatif, keadaan sebuah gen perkembangan yang
dimatikan pada garis keturunan ini tetapi dipakai pada garis lain, sehingga
$B$-lah yang mungkin aktif pada tipe sel lain. $C$: bisu, \omterm{def:b3:chromatin-epigenetics:nucleosome}{heterokromatin}
konstitutif (jalur HP1), lazimnya urutan berulang atau gen yang dibungkam
pada setiap tipe sel.
\end{solution}

\begin{solution}{exo:b3:chromatin-epigenetics:8}
\emph{UBE3A} hanya diekspresikan dari alel ibu di dalam neuron. Penghapusan
yang dibawanya datang dari ayahnya, pada alel yang memang sudah bisu,
sehingga ia sehat; ketika ia menurunkannya, penghapusan itu menjadi
penghapusan dari ibu: tiap anak berpeluang $1/2$ mewarisinya dan lalu
mengidap sindrom Angelman. Andaikan penghapusan itu datang dari ibunya, ia
sendiri akan mengidap sindrom Angelman; pembawa yang sehat hanya mungkin
mewarisinya dari ayah. (Penghapusan seluruh daerah itu dari ayah akan memberi
Prader--Willi; \emph{UBE3A} sendiri tidak cukup untuk itu.)
\end{solution}

\begin{solution}{exo:b3:chromatin-epigenetics:9}
Satu X tanpa \emph{\omterm{def:b3:chromatin-epigenetics:x-inactivation}{Xist}} tidak dapat dinonaktifkan, sehingga X yang lain
dibungkam di setiap sel: penonaktifannya lengkap tetapi tidak lagi acak
(condong), dan embrionya hidup. Kedua X tanpa \emph{\omterm{def:b3:chromatin-epigenetics:x-inactivation}{Xist}}: tanpa
penonaktifan, dua kromosom X aktif, ekspresi terpaut-X dua kali lipat, dan
embrio betina itu mati dini. Sebuah transgen \emph{\omterm{def:b3:chromatin-epigenetics:x-inactivation}{Xist}} pada autosom
menyelimuti autosom tersebut secara cis dan membungkam gennya, sehingga satu
salinan autosom kehilangan fungsi --- percobaan yang menunjukkan
\emph{\omterm{def:b3:chromatin-epigenetics:x-inactivation}{Xist}} memadai untuk pembungkaman secara cis.
\end{solution}

\begin{solution}{exo:b3:chromatin-epigenetics:10}
Sitosina termetilasi terdeaminasi menjadi timina, yang tidak dikenali
permesinan perbaikan, sehingga \omterm{def:b3:chromatin-epigenetics:methylation}{CpG} termetilasi bermutasi menjadi TpG atau CpA
dengan laju tinggi. Hanya urutan yang tak termetilasi di dalam \emph{garis
nutfah}, generasi demi generasi, yang mempertahankan \omterm{def:b3:chromatin-epigenetics:methylation}{CpG-nya}; \omterm{def:b3:chromatin-epigenetics:methylation}{pulau CpG}
mempertahankan miliknya, sehingga pulau itu telah tak termetilasi di dalam
garis nutfah sepanjang hampir seluruh sejarah vertebrata. Gen yang hanya
termetilasi di hati tetap tak termetilasi di dalam sperma dan telur; mutasi
yang timbul di sel hati tidak diwariskan, sehingga metilasi somatik tidak
meninggalkan jejak di dalam urutan.
\end{solution}

\begin{solution}{exo:b3:chromatin-epigenetics:11}
Pada replikasi, \omterm{def:b3:chromatin-epigenetics:nucleosome}{histon} lama yang bertanda dibagi antara kedua duplek anaknya
dan diselingi \omterm{def:b3:chromatin-epigenetics:nucleosome}{histon} baru yang tanpa tanda, sehingga tiap domain anak
bertanda kira-kira separuhnya. HP1 yang terikat pada \omterm{def:b3:chromatin-epigenetics:nucleosome}{nukleosom} bertanda yang
tertahan merekrut SUV39H, yang memetilasi tetangganya yang tanpa tanda;
selama \omterm{def:b3:chromatin-epigenetics:nucleosome}{nukleosom} bertanda cukup rapat untuk membawa seorang \omterm{def:b3:chromatin-epigenetics:histone-code}{penulis} ke tiap
celah, domain itu terisi kembali sebelum pembelahan berikutnya --- umpan
balik yang sama dengan $\mu$ tinggi di dalam teorema. Perambatannya berhenti
pada batas: protein isolator (CTCF), gen yang ditranskripsi kuat dan gen
tRNA, \omterm{def:b3:chromatin-epigenetics:nucleosome}{nukleosom} yang mengusung tanda aktif yang tak serasi (H3K4me3,
\omterm{def:b3:chromatin-epigenetics:histone-code}{asetilasi}) yang tak dapat dipakai \omterm{def:b3:chromatin-epigenetics:histone-code}{penulisnya}, serta titik tempat \omterm{def:b3:chromatin-epigenetics:histone-code}{penghapus}
(demetilase, pergantian HDAC) melampaui \omterm{def:b3:chromatin-epigenetics:histone-code}{penulisnya}. Ujinya: rekrutlah HP1
sesaat ke sebuah pelapor (sebuah fusi dengan domain pengikat DNA yang
pengikatannya dikendalikan obat), lepaskan, lalu ikuti kebisuan pelapor itu
serta H3K9me3-nya sepanjang pembelahan; ramalkan ketahanan yang bergantung
pada \omterm{def:b3:chromatin-epigenetics:histone-code}{penulisnya}, dan kehilangan bila \omterm{def:b3:chromatin-epigenetics:histone-code}{penulis} atau pendimeran \omterm{def:b3:chromatin-epigenetics:histone-code}{pembacanya}
dimutasikan, atau bila sebuah isolator ditempatkan di dalam domain itu.
\end{solution}

\begin{solution}{exo:b3:chromatin-epigenetics:12}
Singkirkanlah: lingkungan bersama (pakan dan kemiskinan cucunya sendiri);
pengaruh rahim --- anak perempuan ibu yang terpapar kelaparan itu masih
janin, dengan sel nutfahnya sudah hadir, sehingga garis nutfah cucunya
terpapar langsung dan pengaruhnya bukan antargenerasi; selisih genetik antara
keluarga yang selamat dan yang tidak; sebab terbalik (massa tubuh itu sendiri
mengubah metilasi di dalam darah); selisih susunan sel darah; dan pengujian
berganda atas banyak \omterm{def:b3:chromatin-epigenetics:methylation}{CpG}. Percobaan pada mencit: paparkan \emph{jantan}
F$_0$ galur murni (tanpa jalur rahim), kawinkan dengan betina yang tak
terpapar atau pakailah pembuahan di luar tubuh serta pemindahan embrio, lalu
ikutilah F$_2$ dan F$_3$ lewat garis ayah sambil mengukur metilasi di dalam sperma
tiap generasi; sebuah tanda yang bertahan pada sperma dan fenotipe F$_3$,
di dalam galur tanpa keragaman genetik, adalah penerusan lewat garis nutfah.
\end{solution}

\begin{solution}{pb:b3:chromatin-epigenetics:1}
\textbf{1.} $6.4\times 10^{9}\times \qty{0.34}{nm} = \qty{2.2}{m}$;
$6.4\times 10^{9}/200 = 3.2\times 10^{7}$ \omterm{def:b3:chromatin-epigenetics:nucleosome}{nukleosom}.
\textbf{2.} \omterm{def:b3:chromatin-epigenetics:nucleosome}{Histon}: $3.2\times 10^{7}\times 1.08\times 10^{5} =
3.5\times 10^{12}$ Da $= \qty{5.7}{pg}$. DNA: $6.4\times 10^{9}\times
650 = 4.2\times 10^{12}$ Da $= \qty{6.9}{pg}$. Massanya sebanding.
\textbf{3.} Nukleus: $\tfrac{4}{3}\pi\,(\qty{5}{\micro m})^{3} =
\qty{520}{\micro m^3}$. Inti: $\pi\,(5.5)^{2}\times 5.5 =
\qty{520}{nm^3} = 5.2\times 10^{-7}\,\unit{\micro m^3}$; totalnya
$3.2\times 10^{7}\times 5.2\times 10^{-7} = \qty{17}{\micro m^3}$:
sekitar \qty{3}{\%} nukleus.
\textbf{4.} $3.2\times 10^{7}\times \qty{11}{nm} = \qty{0.35}{m}$;
\omterm{prop:b3:chromatin-epigenetics:compaction}{rasio pemampatannya} $2.2/0.35 \approx 6$.
\textbf{5.} $6.4\times 10^{9}/2\times 10^{5} = \num{32000}$ simpal, dengan
\num{1000} \omterm{def:b3:chromatin-epigenetics:nucleosome}{nukleosom} pada tiap simpal.
\textbf{6.} $0.21^{2}\times 6.4\times 10^{9} = 2.8\times 10^{8}$;
teramati $5.6\times 10^{7}$, jadi \qty{80}{\%} hilang, lewat deaminasi
5-metilsitosina menjadi timina (transisi C$\to$T yang tak terperbaiki pada
\omterm{def:b3:chromatin-epigenetics:methylation}{CpG} termetilasi).
\textbf{7.} $p_{n+1} = 0.99\,p_{n} + 0.02\,(1 - p_{n}) = 0.97\,p_{n} +
0.02$; $p^{*} = 0.02/0.03 = 0.67$.
\textbf{8.} $p_{n} = 0.67 + 0.33\times 0.97^{n}$: $n = 10$: $0.91$;
$n = 50$: $0.74$; $n = 200$: $0.67$ (kelebihannya $7\times 10^{-4}$).
\textbf{9.} $\ln 2/(-\ln 0.97) = 22.8 \approx 23$ pembelahan; pada satu
pembelahan seminggu, sekitar \num{23} minggu --- lima bulan.
\textbf{10.} Tiap replikasi memasangkan satu untai lama, yang masih
termetilasi, dengan satu untai baru yang tak dimetilasi enzim mana pun;
untai lamanya tidak bertambah maupun berkurang, cacah untainya berlipat dua,
sehingga fraksi termetilasinya separuh. $0.80/2^{n} < 0.05$ menuntut
$2^{n} > 16$: lima pembelahan ($0.025$; empat pembelahan memberi tepat
\qty{5}{\%}).
\textbf{11.} Rasionya $0.9995 - 0.0005 = 0.999$; waktu paruhnya $\ln
2/(-\ln 0.999) \approx 690$ pembelahan, sekitar \num{13} tahun pada satu
pembelahan seminggu.
\textbf{12.} Gen yang bisu selama seumur hidup, yakni sekitar \num{4000}
pembelahan mingguan, menuntut waktu paruh ratusan pembelahan, dan hanya
umpan balik pertanyaan~11 yang memberinya: \omterm{def:b3:chromatin-epigenetics:histone-code}{pembaca} tandanya (MeCP2, HP1)
merekrut \omterm{def:b3:chromatin-epigenetics:histone-code}{penulisnya} (metilase H3K9, DNMT3), sehingga domain tanda yang rapat
menyalin dirinya dengan daya guna yang tak dimiliki enzim tunggal mana pun.
\textbf{13.} Satu X mengusung alel jingga, X yang lain alel hitam; tiap sel
kulit telah membungkam satu X dan hanya mengekspresikan yang lain;
keturunan sebuah sel mempertahankan pilihannya, sehingga satu bercak adalah
satu klon (atau sekelompok klon bersebelahan) yang membuat pilihan yang sama.
\textbf{14.} Seluruh $N$ sel memilih X pembawa jingga, atau seluruhnya
memilih X yang lain: $2\times (1/2)^{N} = 2^{1-N}$.
\textbf{15.} $2^{1-N} = 10^{-9}$: $N - 1 = \log_{2} 10^{9} = 29.9$,
$N \approx 30$ sel.
\textbf{16.} $\qty{3000}{cm^2}/\qty{15}{cm^2} = 200$ bercak, jauh lebih
banyak daripada $30$ pendiri: selama pertumbuhan keturunan seorang pendiri
terpencar oleh percampuran dan perpindahan sel, sehingga satu klon membentuk
beberapa pulau (dan tetangga yang berpilihan sama menyatu, yang bekerja ke
arah sebaliknya tetapi lebih lemah).
\textbf{17.} Satu \omterm{def:b3:chromatin-epigenetics:x-inactivation}{badan Barr} tiap X di luar X yang pertama: XXY satu, XXX
dua, XY tidak ada.
\textbf{18.} Dapat ditoleransi bila gennya memiliki homolog Y yang
berfungsi, sehingga jantan pun memiliki dua salinan (gen pseudoautosom dan
pasangan seperti \emph{KDM5C}/\emph{KDM5D}); hal itu penting pada aneuploidi
X --- pada sindrom Turner (XO) gen yang lolos hanya hadir dalam satu salinan,
dan itulah sebabnya individu XO bukan sekadar perempuan normal
(haploinsufisiensi \emph{SHOX} dan perawakan pendek), sedangkan pada XXY gen
yang lolos justru berlebih dosis.
\textbf{19.} $1/\num{15000} = 6.7\times 10^{-5}$; penghapusannya $0.70\times
6.7\times 10^{-5} = 4.7\times 10^{-5}$ (satu dari \num{21000}); disomi
ibunya $1.7\times 10^{-5}$ (satu dari \num{60000}).
\textbf{20.} Angelman adalah kehilangan satu gen, \emph{UBE3A}, sehingga
satu mutasi titik pada alel ibu sudah memadai. Prader--Willi adalah
kehilangan beberapa gen yang diekspresikan dari ayah sekaligus
(\emph{SNRPN}, kelompok RNA kecilnya); mutasi titik menonaktifkan salah satu
saja dan memberi paling banter fenotipe sebagian, tak pernah sindrom
seutuhnya.
\textbf{21.} \qty{20}{\%} anaknya menerima penghapusan dari ayah dan
mengidap Prader--Willi; tidak ada yang mengidap Angelman, yang menuntut
penghapusan dari ibu.
\textbf{22.} Tiga kromosom, dua dari ibu dan satu dari ayah, satu hilang
secara acak: yang dari ayah hilang dengan peluang $1/3$, meninggalkan disomi
ibu dan Prader--Willi; dengan peluang $2/3$ satu salinan dari ibu yang
hilang dan anaknya normal.
\textbf{23.} Gen dari ayah semestinya memacu tuntutan pada ibu;
kehilangannya (Prader--Willi) semestinya menurunkan tuntutan --- dan itu
cocok dengan lemahnya isapan dan gagal tumbuh bayi baru lahir --- sedangkan
nafsu makan rakus yang datang kemudian justru kebalikannya dan lebih sukar
dicocokkan (satu bacaan yang diusulkan ialah bahwa gen dari ayah biasanya
mengarahkan anak menuju penyapihan). Angelman, kehilangan sebuah gen dari
ibu, semestinya menaikkan tuntutan; isapan yang berkepanjangan dan seringnya
terbangun malam pada anak Angelman memang cocok, sedangkan kejang dan
hilangnya bicara sama sekali bukan soal sumber daya.
\textbf{24.} Sasarlah transkrip antisens dari ayah
(\emph{UBE3A-ATS}) dengan oligonukleotida antisens yang mengurainya
atau menghalangi pemanjangannya: \emph{UBE3A} dari ayah masih utuh dan
dibungkam hanya secara cis oleh RNA itu, sehingga membuang RNA-nya memulihkan
ekspresinya. Metilasi \omterm{def:b3:chromatin-epigenetics:imprinting}{daerah kendali pencetakannya} tidak tersentuh,
sehingga \emph{SNRPN} dan gen tercetak yang lain mempertahankan pola
tetuanya yang normal.
\textbf{25.} Waktu paruh sebuah tanda: sekitar $23$ pembelahan tanpa umpan
balik, sekitar $690$ dengan umpan balik; bulu belang tiga diputuskan oleh
populasi pendiri sekitar $30$ sel.
\end{solution}
